\scanpage{260}
\uline{Corollaire} 1.8. Soit $f:X\to Y$ un morphisme propre d'espaces analytiques \inserted{sur $\mathbb C$}, et $y\in Y$. Conditions équivalentes :

(i) $\exists$ un voisinage ouvert $U$ de $y$ tel que $X|U$ soit \uline{proj} sur $U$ (cf Sém. Cartan).

(ii) Pour tout entier $n$, $X_n=X\times_Y Y_n$ (i.e. $Y_n=\operatorname{Specan}\mathcal O_y/\mathfrak m_y^{n+1}$) est \deleted{\unsure{proj}} un espace analytique projectif.

\uline{Dém}. Réf : Sém. Cartan ....

Exemple

\margin{\figure{\begin{tikzpicture}[x=1cm,y=1cm]\draw (0,4.5)--(.06,3.1)--(.1,1.4)..controls(.13,.5)and(.3,.3)..(.18,0);\draw (.06,4.45)--(.15,3.2);\end{tikzpicture}}}
\uline{Corollaire} \deleted{\unsure{1.8|1.9}} 1.9. \deleted{Supposons} Soit $f:X\unsure{-|\to}Y$ un morphisme propre d'espaces analytiques sur $\mathbb C$, \deleted{\unsure{et} tel que \unsure{dim|les} $\unsure{X_y}$} dont les fibres soient de dim $\leq 1$. Alors $f$ est localement projectif.

Résulte du

\uline{Lemme} 1.9. Un espace analytique compact de dim $\leq 1$ est projectif.

\mspage{259}
